\documentclass{article}
\usepackage[T1]{fontenc}
\usepackage{lmodern}
\usepackage{graphicx}
\usepackage{amsmath,amssymb}
\usepackage[a4paper,margin=2cm]{geometry}
\usepackage[absolute,overlay]{textpos}
\setlength{\TPHorizModule}{1mm}
\setlength{\TPVertModule}{1mm}
\usepackage[indonesian]{babel}
\usepackage{hyperref}
\setlength{\emergencystretch}{2em}
% unit: Halaman 36

\begin{document}

% === Halaman 36 (python/native) ===

36

• Enkripsi dilakukan dengan mencari titik potong huruf plainteks dengan 

huruf kunci:


Plainteks 

















K 

U 

N 

C 

I 


A 

B 

C 

D 

E 

F 

G 

H 

I 

J 

K 

L 

M 

N 

O 

P 

Q 

R 

S 

T 

U 

V 

W 

X 

Y 

Z 

a 

A 

B 

C 

D 

E 

F 

G 

H 

I 

J 

K 

L 

M 

N 

O 

P 

Q 

R 

S 

T 

U 

V 

W 

X 

Y 

Z 

b 

B 

C 

D 

E 

F 

G 

H 

I 

J 

K 

L 

M 

N 

O 

P 

Q 

R 

S 

T 

U 

V 

W 

X 

Y 

Z 

A 

c 

C 

D 

E 

F 

G 

H 

I 

J 

K 

L 

M 

N 

O 

P 

Q 

R 

S 

T 

U 

V 

W 

X 

Y 

Z 

A 

B 

d 

D 

E 

F 

G 

H 

I 

J 

K 

L 

M 

N 

O 

P 

Q 

R 

S 

T 

U 

V 

W 

X 

Y 

Z 

A 

B 

C 

e 

E 

F 

G 

H 

I 

J 

K 

L 

M 

N 

O 

P 

Q 

R 

S 

T 

U 

V 

W 

X 

Y 

Z 

A 

B 

C 

D 

f 

F 

G 

H 

I 

J 

K 

L 

M 

N 

O 

P 

Q 

R 

S 

T 

U 

V 

W 

X 

Y 

Z 

A 

B 

C 

D 

E 

g 

G 

H 

I 

J 

K 

L 

M 

N 

O 

P 

Q 

R 

S 

T 

U 

V 

W 

X 

Y 

Z 

A 

B 

C 

D 

E 

F 

h 

H 

I 

J 

K 

L 

M 

N 

O 

P 

Q 

R 

S 

T 

U 

V 

W 

X 

Y 

Z 

A 

B 

C 

D 

E 

F 

G 

i 

I 

J 

K 

L 

M 

N 

O 

P 

Q 

R 

S 

T 

U 

V 

W 

X 

Y 

Z 

A 

B 

C 

D 

E 

F 

G 

H 

j 

J 

K 

L 

M 

N 

O 

P 

Q 

R 

S 

T 

U 

V 

W 

X 

Y 

Z 

A 

B 

C 

D 

E 

F 

G 

H 

I 

K 

K 

L 

M 

N 

O 

P 

Q 

R 

S 

T 

U 

V 

W 

X 

Y 

Z 

A 

B 

C 

D 

E 

F 

G 

H 

I 

J 

l 

L 

M 

N 

O 

P 

Q 

R 

S 

T 

U 

V 

W 

X 

Y 

Z 

A 

B 

C 

D 

E 

F 

G 

H 

I 

J 

K 

m 

M 

N 

O 

P 

Q 

R 

S 

T 

U 

V 

W 

X 

Y 

Z 

A 

B 

C 

D 

E 

F 

G 

H 

I 

J 

K 

L 

n 

N 

O 

P 

Q 

R 

S 

T 

U 

V 

W 

X 

Y 

Z 

A 

B 

C 

D 

E 

F 

G 

H 

I 

J 

K 

L 

M 

o 

O 

P 

Q 

R 

S 

T 

U 

V 

W 

X 

Y 

Z 

A 

B 

C 

D 

E 

F 

G 

H 

I 

J 

K 

L 

M 

N 

p 

P 

Q 

R 

S 

T 

U 

V 

W 

X 

Y 

Z 

A 

B 

C 

D 

E 

F 

G 

H 

I 

J 

K 

L 

M 

N 

O 

q 

Q 

R 

S 

T 

U 

V 

W 

X 

Y 

Z 

A 

B 

C 

D 

E 

F 

G 

H 

I 

J 

K 

L 

M 

N 

O 

P 

r 

R 

S 

T 

U 

V 

W 

X 

Y 

Z 

A 

B 

C 

D 

E 

F 

G 

H 

I 

J 

K 

L 

M 

N 

O 

P 

Q 

s 

S 

T 

U 

V 

W 

X 

Y 

Z 

A 

B 

C 

D 

E 

F 

G 

H 

I 

J 

K 

L 

M 

N 

O 

P 

Q 

R 

t 

T 

U 

V 

W 

X 

Y 

Z 

A 

B 

C 

D 

E 

F 

G 

H 

I 

J 

K 

L 

M 

N 

O 

P 

Q 

R 

S 

u 

U 

V 

W 

X 

Y 

Z 

A 

B 

C 

D 

E 

F 

G 

H 

I 

J 

K 

L 

M 

N 

O 

P 

Q 

R 

S 

T 

v 

V 

W 

X 

Y 

Z 

A 

B 

C 

D 

E 

F 

G 

H 

I 

J 

K 

L 

M 

N 

O 

P 

Q 

R 

S 

T 

U 

w 

W 

X 

Y 

Z 

A 

B 

C 

D 

E 

F 

G 

H 

I 

J 

K 

L 

M 

N 

O 

P 

Q 

R 

S 

T 

U 

V 

x 

X 

Y 

Z 

A 

B 

C 

D 

E 

F 

G 

H 

I 

J 

K 

L 

M 

N 

O 

P 

Q 

R 

S 

T 

U 

V 

W 

y 

Y 

Z 

A 

B 

C 

D 

E 

F 

G 

H 

I 

J 

K 

L 

M 

N 

O 

P 

Q 

R 

S 

T 

U 

V 

W 

X 

z 

Z 

A 

B 

C 

D 

E 

F 

G 

H 

I 

J 

K 

L 

M 

N 

O 

P 

Q 

R 

S 

T 

U 

V 

W 

X 

Y 

Gambar 4.3  Enkripsi huruf T dengan kunci s 

Plainteks  : terimapesanangulaikambing

Kunci 

 : soreangsoreangsoreangsore

Cipherteks: LSIMMNVWGRRAAMMZRMKNSTWEK

\end{document}
